\begin{figure}[t]
\begin{minipage}[b]{0.44\linewidth}
  \centering
  \small
  \setlength{\tabcolsep}{4pt}
  \begin{tabular}{l c}
  \toprule
  \textbf{Condition group} & \textbf{Mean IF} \\
  \midrule
  \emph{Neutral}       & $31$~\textit{[10--52]}\% \\
  \emph{Classify}      & $38$~\textit{[21--56]}\% \\
  \emph{Random-facts}  & $50$~\textit{[35--66]}\% \\
  \emph{Task-based}    & $43$~\textit{[31--56]}\% \\
  \bottomrule
  \end{tabular}
  \captionof{table}{Mean instruction-following across the 13 models for the $2\times2$ engagement~$\times$~diversity follow-up ($T{=}0$, no-hint, 35 trials/cell; 95\% CI across models in brackets). This suggests that output diversity and length, not question engagement, accounts for the \emph{task-based} robustness gap.}
  \label{tab:followup-main}
\end{minipage}%
\hfill
\begin{minipage}[b]{0.52\linewidth}
  \centering
  \includegraphics[width=\linewidth]{\plotsroot/dynamic/f4_followup_curves.png}
  \caption{Mean IF rate vs.\ $N$ by condition group, averaged over the 13 models (bands: $\pm1$ SE across models). \emph{random-facts} (output diversity) tracks the \emph{task-based} curve, whereas \emph{classify} (question engagement) stays near the \emph{neutral}, \emph{fixed-output} condition baseline.}
  \label{fig:followup-curves}
\end{minipage}
\end{figure}
